\section{武周的区域治理、宗教与边防} \label{sec:wu-zhou-regional-governance}  oindent	extbf{材料的入口。} 武周的叙事不能只从皇帝、国号或一场战役的结局开始。临朝、徐敬业起兵、建周、安西恢复、营州契丹起事与复唐这些节点把中枢命令放进了州县、军镇、仓储、道路、寺院和家庭能够实际接触的场所。账本所锁定的事件编号为 ST-683-GAOZONG-DEATH、ST-684-LI-JINGYE、ST-690-ZHOU-FOUNDATION、ST-692-ANXI-RECOVERY、ST-695-KHITAN-REBELLION、ST-705-TANG-RESTORATION；它们提供可复核的年序和行动者，却不自动提供所有地区的同一种经验。国家文书最擅长保存颁令、任免和胜败，地方社会往往只在逃亡、诉讼、征发、迁徙或墓志的一行文字里短暂出现。因而，叙事必须将“朝廷说要如何办”与“某一处实际如何承受”分开，并承认后者常常无从完整恢复。   oindent	extbf{财政与行政的尺度。} 每一次改制、出兵或册封，最终都要落到粮食如何集中、差役由谁承担、船只沿何处行驶、军队在哪一座城停留的问题。户籍、手实、仓簿和契约让研究者看见国家分类的技术；它们却多保存于敦煌、吐鲁番或个别精英家庭，不能由一个绿洲、一个县或一座墓葬推出全帝国的税率、人口和服役结构。反过来，正史食货志、地理志与会要能够给出制度的广度，却常把缓慢、参差、反复的执行压缩成一条整齐沿革。将二者并读，所能得到的不是一幅完整账本，而是对地区差异更严格的提问：哪些收入进入中枢，哪些留在军镇；哪些家庭进入名册，哪些人通过迁徙、依附或隐匿逃离分类。   oindent	extbf{区域与边地。} 边疆并非中原国家向外铺开的空白地图。山地、草原、绿洲、海港和河谷各有不同的补给条件与政治语言；“归附”“册封”“置州”或“收复”首先是材料中可确认的行政动作，不能直接换算为稳定的统治深度。军事传记的兵数、歼获与行军路线常服务军功和合法性，邻国史书、碑铭和多语文书又有不同的日期与目的。写作时应把这些材料的错位保存下来：中原本纪可确定某道诏令或某次撤军的年序，地方遗存和外部记录则提醒我们不要让朝廷的分类成为行动者唯一的自称。   oindent	extbf{宗教、文化与社会关系。} 寺院、译场、学校、科场、市场和家族祭祀并不是政治史以外的装饰。它们使文字、财产、名望与人能够跨区域移动，也使国家在征用、保护、禁毁和授官时进入普通人可感的世界。僧传、诗文、碑志与笔记尤其需要谨慎：它们保存了制度感受和地方记忆，却带着典范、修辞和传抄层次，不能把一篇表文或一方功德碑变成普遍信仰的证明。文化史的价值在于打断单一的官僚叙事，而不是用另一套同样整齐的印象取代它。   oindent	extbf{史料限度与叙事判断。} 本节主要依托两《唐书》、通鉴、诏令会要、吐蕃与突厥材料、碑刻。这些材料各有时间距离和观察位置：官修史书会以继承与治乱组织事件，制度汇编会将实践化为条目，出土材料只照亮保存下来的局部。故而可以确认某项政策被颁行、某座城市易手、某位首领获授名号；不能据此断定每个乡村已经服从、每一笔收入都已入库，或每个行动者怀有同一动机。保持这种限度，并不会削弱叙事，反而使国家能力、地方策略和社会代价重新进入同一时间框架。  